%%% FP2001.tex 05.10.2000 - Vesa Kolhinen <vesa.kolhinen@phys.jyu.fi>
%%% Original version:
%%% FP2000.tex 12.10.1999 - Eero Noponen <eero.noponen@hut.fi>
%%%
%%%
%%%
%%% Sample LaTeX file for the abstracts for the Proceedings of the
%%% Annual Meeting of the Finnish Physical Society.

\documentclass[12pt]{article}
\usepackage{fp2009}
\usepackage[T1]{fontenc}
\usepackage[latin1]{inputenc}
\usepackage{hyperref}

%%% Remove the following line if Times font is unavailable.
\usepackage{times}

%%% The following line is required if the document contains figures.
\usepackage{graphicx}

%%% The following line is for LaTeX 2.09.
%\documentstyle[FP2001,12pt]{article}

\begin{document}

%%%
%%% The title of the abstract is given here as \title{TITLE TEXT}.
%%%
\title{USING GIT DISTRIBUTED VERSION CONTROL TOOL TO MAINTAIN PHYSICS
  CODE AND INTEROPERATE WITH CERN SUBVERSION REPOSITORIES}

%%%
%%% The list of authors; underline the name of the author who is responsible
%%% for the presentation of the paper.
%%%
A. Heikkinen \underline{P. Kaitaniemi}

%%%
%%% The address.
%%%
P.O.Box 64 (Gustaf H\"allstr\"omin katu 2), FIN-00014 University of Helsinki, Finland\\
email: pekka.kaitaniemi@gmail.com

%%% Do not remove.
\vspace{\baselineskip}

%%%
%%% Body text begins here.
%%%

{\tt Git} is a \emph{distributed version control} tool developed by
Linus Torvalds and the Linux kernel development team.

Version control allows developers to keep track of project history in
a systematic and detailed manner. The project can be a data analysis
or simulation program source code, or even LaTeX code of a paper or
thesis. In addition to tracking history, version control tools allow
one to merge contributions between several authors.

%%% The following minipage structures are needed to position the figure
%%% beside the text paragraph.
%%% For a standard figure environment, see the example below.

\begin{minipage}{9cm}
Figures may be included in the submitted files. Here we show how a
figure can be included in a \LaTeX\ file.
\end{minipage}
\hfill
\begin{minipage}{5cm}
\includegraphics[width=4cm]{logo.eps}
\end{minipage}

%%%
%%% A more typical Figure.
%%%
%\begin{figure}[!h]
%\centering\includegraphics[width=4cm]{logo.eps}
%\caption{Mallikuva.}
%\end{figure}

%%%
%%% The bibliography: the references are listed here.
%%%
\begin{thebibliography}{9}

\bibitem{ref1}
A. Fyysikko and B. Fyysikko,
\href{http://link.aps.org/abstract/PR/v108/p1175}{Fysiikan Lehti 108
(1957) 1175}.
\end{thebibliography}

\end{document}
